\documentclass{article}
\usepackage[utf8]{inputenc}
\usepackage[T2A]{fontenc}
\usepackage[english,russian]{babel}

\title{Пример рукописи для проверки гейтов}
\author{Автор}

\begin{document}
\maketitle

\section{Введение}
Настоящая работа посвящена воспроизводимости вычислительных экспериментов.
Мы опираемся на открытые данные, опубликованные ранее \cite{smith2020},
и сопоставляем их с собственными измерениями \cite{ivanov2024}.

\section{Методы}
Выборка включала 42 участника. Критерии включения описаны в реестре протокола.
Обработка данных выполнялась скриптом \texttt{analysis/run\_all.sh}, версии
фиксируются в \texttt{analysis/environment.lock}.

\section{Результаты}
Разница между группами составила 3,1 балла по шкале HAM-D.
Все таблицы пересчитываются из сырых данных одной командой.

\bibliographystyle{plain}
\bibliography{refs}
\end{document}
